\documentclass[10pt]{article}
\usepackage[a4paper,margin=0.72in]{geometry}
\usepackage{amsmath,amssymb}
\usepackage{booktabs}
\usepackage{graphicx}
\usepackage[hidelinks]{hyperref}
\usepackage{microtype}
\usepackage{enumitem}
\usepackage{array}

\setlist{nosep,leftmargin=*}
\setlength{\parindent}{0pt}
\setlength{\parskip}{4pt}
\setlength{\textfloatsep}{7pt plus 2pt minus 2pt}
\setlength{\floatsep}{6pt plus 2pt minus 2pt}
\setlength{\intextsep}{6pt plus 2pt minus 2pt}
\newcolumntype{L}[1]{>{\raggedright\arraybackslash}p{#1}}

\title{\vspace{-1.2cm}K=2 Lazy Ring with One Directed Shortcut: Beta Sweep for the Double Peak}
\author{}
\date{}

\begin{document}
\maketitle
\vspace{-1.2cm}

\section*{Conclusion}
For the tested case
$N=100$, $q=2/3$, $u=6$, $v=56$, and $n_0=1,\ldots,6$, the proposed boundary is
\[
  \beta_c=\frac{2q}{(1-q)N}=0.04.
\]
I use R3 as the main criterion for a visually resolved double peak.  Let
$t_1<t_2$ be the first two filtered local-maximum times, with
$h_1=f(t_1)$ and $h_2=f(t_2)$, and let $h_v=f(t_v)$ be the minimum value of
$f(t)$ between those two peaks.  Define
\[
  \rho=\frac{h_v}{\min(h_1,h_2)}.
\]

The criteria used below are:
\[
\begin{array}{ll}
\mathrm{R1}: & \text{two filtered local maxima, separated in time and with balanced heights},\\
\mathrm{R2}: & \mathrm{R1}+h_v/\max(h_1,h_2)\leq0.9,\\
\mathrm{R3}: & \mathrm{R1}+\rho=h_v/\min(h_1,h_2)\leq0.8,\\
\mathrm{R4}: & \mathrm{R1}+\rho=h_v/\min(h_1,h_2)\leq0.6.
\end{array}
\]
R2 is a permissive ``there is a dip'' diagnostic; R3 is the main visual
double-peak criterion because it requires the valley to be clearly below the
smaller peak.

The assumption behind the proposed boundary is that the transition can be
captured by an $n_0$-independent necessary condition in $\beta$.  The sweep
supports this assumption only in the visual sense: under R3, no clear double
peak remains at or above $\beta_c$.  The reference case $n_0=1$,
$\beta=0.01$ does give a clear double peak, but at $\beta=0.04$ the observed
values are already $\rho\simeq0.887$--$0.891$ for $n_0=1,2,3$, so the feature
is better described as a shallow shoulder or marginal second peak.

The proposed boundary is therefore supported as a cutoff for a visually clear
double peak, but not as a strict cutoff for the existence of two mathematical
local maxima.  Under the weaker R1/R2 rules, two local maxima can still be
detected above $\beta_c$ for $n_0=1,2,3$, up to about
$\beta=0.060$--$0.062$.  The starting point $n_0$ matters, so the expression
cannot be the whole condition.

\section*{Question}
I test the proposed condition for the first-passage distribution on a lazy
nearest-neighbour ring with one directed shortcut.  The model is
\[
  N=100,\qquad K=2,\qquad q=\frac23,\qquad u=6,\qquad v=56,
\]
with paper-indexed sites.  At the shortcut source, a fraction of the lazy
self-loop probability is moved to the directed shortcut:
\[
  P(u\to v)=\beta(1-q), \qquad
  P(u\to u)=(1-\beta)(1-q).
\]
The sweep uses start positions $n_0=1,\ldots,6$ and beta values from $0$ to
$0.08$, with extra points at the reference values.  The default target is the
named endpoint $v=56$; the script can be rerun with another target convention
using \texttt{--target-paper}.

The proposed boundary is
\[
  \beta_c=\frac{2q}{(1-q)N}=0.04.
\]

\begin{figure}[!htbp]
  \centering
  \includegraphics[width=0.55\linewidth]{../artifacts/figures/luca_k2_fixed_shortcut_ring_geometry.pdf}
  \caption{Concrete geometry of the tested ring.  The grey circle is the
  $N=100$, $K=2$ nearest-neighbour ring; yellow points show the swept starts
  $n_0=1,\ldots,6$; the red curved arrow is the directed shortcut
  $u=6\to v=56$, whose endpoint is also the absorbing target.}
\end{figure}

\section*{Framework}
For each $(n_0,\beta)$ I compute $f(t)$ up to
$t_{\max}=12000$.  Because the phrase ``double peak'' can mean either
``two mathematical local maxima'' or ``two peaks clearly visible by eye'', I
record several definitions.  The visual definitions use
$h_v/\min(h_1,h_2)$ because a curve only looks like two peaks if the valley is
below both peaks, especially the smaller one.

The quantities used in the rules are:
\[
  t_1<t_2,\qquad h_1=f(t_1),\qquad h_2=f(t_2),\qquad
  t_v=\arg\min_{t_1<t<t_2} f(t),\qquad h_v=f(t_v).
\]
Here $t_1$ and $t_2$ are the first two filtered local-maximum times, $t_v$ is
the intervening valley time, $h_2/h_1$ measures the relative peak heights,
$h_v/\max(h_1,h_2)$ measures whether there is any valley below the larger peak,
and $h_v/\min(h_1,h_2)$ measures whether the valley is visibly below both
peaks.

The operational rules are:
\begin{description}[style=nextline,leftmargin=1.3cm,labelwidth=1.1cm]
  \item[R1] \textbf{Two local maxima}: two filtered local peaks with
  $t_2-t_1\geq10$ and balanced peak heights
  $0.1\leq h_2/h_1\leq10$.
  \item[R2] \textbf{Permissive valley}: R1 plus
  $h_v/\max(h_1,h_2)\leq0.9$.
  \item[R3] \textbf{Visual double peak}: R1 plus
  $\rho=h_v/\min(h_1,h_2)\leq0.8$.
  \item[R4] \textbf{Strong visual double peak}: R1 plus
  $\rho=h_v/\min(h_1,h_2)\leq0.6$.
\end{description}

The difference between R2 and R3 is intentional.  R2 uses
$h_v/\max(h_1,h_2)$, so it only asks whether the curve has a dip below the
larger of the two peaks.  This is a permissive diagnostic: it confirms that a
valley exists, but it can still count cases where the valley is almost as high
as the smaller peak.  Such cases may have two mathematical local maxima without
looking like two clearly separated peaks.

R3 uses $h_v/\min(h_1,h_2)$ instead.  This is the more relevant visual test
because a double peak is only clearly resolved if the valley is below both
peaks, especially below the smaller one.  The cutoff $\rho\leq0.8$ is therefore
used as the main ``visible double peak'' rule: it requires a valley at least
about 20\% below the smaller peak.  R4 keeps the same logic but uses the
stricter cutoff $\rho\leq0.6$ to identify only strongly separated cases.

This threshold should be read as a practical visual convention rather than an
absolute mathematical constant.  If the visual rule were relaxed toward
$\rho\leq0.9$, the boundary case $\beta=\beta_c$ would become borderline,
because the measured values there are $\rho\simeq0.887$--$0.891$ for
$n_0=1,2,3$.  If the rule is tightened to R4, the detected double-peak region
shrinks to smaller beta values.  This sensitivity is why R1/R2, R3, and R4 are
reported separately.

\section*{Main Readout}
\begin{table}[!htbp]
  \centering
  \caption{Summary of the fixed-shortcut sweep.  ``Double at/above
  $\beta_c$'' counts grid points satisfying the R3 visual double-peak rule and
  $\beta\geq0.04$.}
  \begin{tabular}{rrrrrr}
\toprule
$n_0$ & target & $\beta_c$ & any double & max $\beta$ double & double at/above $\beta_c$ \\
\midrule
1 & 56 & 0.040 & yes & 0.030 & 0 \\
2 & 56 & 0.040 & yes & 0.030 & 0 \\
3 & 56 & 0.040 & yes & 0.030 & 0 \\
4 & 56 & 0.040 & no & -- & 0 \\
5 & 56 & 0.040 & no & -- & 0 \\
6 & 56 & 0.040 & no & -- & 0 \\
\bottomrule
\end{tabular}

\end{table}

For the reference point mentioned in the request, $n_0=1$ and $\beta=0.01$,
the sweep gives a clear double peak with
\[
  t_1=37,\qquad t_v=226,\qquad t_2=1077,
\]
and peak-height ratio $h_2/h_1\approx1.73$.

\section*{Observed Phenomenology}
\begin{table}[!htbp]
  \centering
  \footnotesize
  \caption{Observed behaviour under different double-peak definitions.
  Here $\rho_c$ is $\rho=h_v/\min(h_1,h_2)$ at $\beta_c=0.04$ for
  $n_0=1,2,3$.  The range $0.887$--$0.891$ means that the valley is already
  close to the smaller peak at the proposed boundary.}
  \begin{tabular}{L{0.21\linewidth} c L{0.39\linewidth} c}
\toprule
Definition & Cases with $\beta\geq\beta_c$ & Maximum passing $\beta$ & $\rho_c$ \\
\midrule
two local maxima & 34 & n0=1: 0.062; n0=2: 0.060; n0=3: 0.060; n0=4: none; n0=5: none; n0=6: none & 0.887-0.891 \\
permissive valley & 34 & n0=1: 0.062; n0=2: 0.060; n0=3: 0.060; n0=4: none; n0=5: none; n0=6: none & 0.887-0.891 \\
visual double peak & 0 & n0=1: 0.030; n0=2: 0.030; n0=3: 0.030 & 0.887-0.891 \\
strong visual double peak & 0 & n0=1: 0.018; n0=2: 0.018; n0=3: 0.018 & 0.887-0.891 \\
\bottomrule
\end{tabular}

\end{table}

The different rules separate two effects.  Under local-maximum rules, two
maxima can still be detected above $\beta_c$ for $n_0=1,2,3$.  Under the visual
rules based on $h_v/\min(h_1,h_2)$, no clear double peak appears at or above
$\beta_c$.  At $\beta=0.04$, the detector can still find two local maxima for
$n_0=1,2,3$, but $h_v/\min(h_1,h_2)\approx0.89$, so the valley is shallow and
the profile is visually closer to a shoulder or marginal second peak than to a
clearly separated double peak.

\clearpage
\section*{Visual Evidence}
\begin{figure}[!htbp]
  \centering
  \includegraphics[width=0.88\linewidth]{../artifacts/figures/luca_k2_fixed_shortcut_annotated_double_peak.pdf}
  \caption{The concrete double-peak curve for the reference case
  $n_0=1$, $\beta=0.01$.  The first peak, valley, and second peak are marked
  explicitly.}
\end{figure}

\begin{figure}[!htbp]
  \centering
  \includegraphics[width=0.92\linewidth]{../artifacts/figures/luca_k2_fixed_shortcut_double_peak_map.pdf}
  \caption{R3 visual double-peak flag over the $(n_0,\beta)$ grid.  The dashed
  red line is the proposed $\beta_c=0.04$.}
\end{figure}

\begin{figure}[!htbp]
  \centering
  \includegraphics[width=0.92\linewidth]{../artifacts/figures/luca_k2_fixed_shortcut_peak_times.pdf}
  \caption{Detected peak locations across beta.  The late peak remains
  separated from the early peak for $n_0=1,2,3$ beyond $\beta_c$ on this grid.}
\end{figure}

\begin{figure}[!htbp]
  \centering
  \includegraphics[width=0.92\linewidth]{../artifacts/figures/luca_k2_fixed_shortcut_representative_fpt.pdf}
  \caption{Representative first-passage curve shapes for selected beta values.
  Each curve is normalized by its own maximum to make the peak and valley
  geometry visible on a linear scale.}
\end{figure}

\clearpage
\section*{Interpretation}
Under the literal target convention $target=v=56$, the sweep does not support
the strong statement that $\beta>\beta_c$ rules out \emph{any two local
maxima}.  In the default grid, local maxima can persist up to
$\beta\approx0.062$ for $n_0=1$ and up to $\beta\approx0.060$ for $n_0=2,3$.
However, under the more appropriate visual criterion
$h_v/\min(h_1,h_2)\leq0.8$, no clear double peak survives at or above
$\beta_c=0.04$.

This agrees with the caveat that the expression cannot be the whole story:
the starting point $n_0$ materially changes the outcome, so any complete
condition must include the geometry of the start, shortcut source, shortcut
destination, and target.

\section*{Reproducibility}
Run from the repository root:
\begin{verbatim}
.venv/bin/python research/reports/ring_lazy_jump_ext_rev2/code/\
luca_k2_fixed_shortcut_sweep.py
\end{verbatim}
The run writes the config, metrics, summary table, selected curves, and all
figures under this report's \texttt{artifacts/} directory.

\end{document}
